\scanpage{328}
\mspage{29}

je n'ai eu connaissance de l'Enterrement qu'en avril 1984, au bout de 14 ans...) Et rares ont été ceux parmi mes anciens amis qui, ayant senti ces effluves, ne s'en soient réjouis en leur for intérieur et n'aient fait chorus à la dérision dont j'étais l'objet. Ce faisant, ils croyaient qu'à toutes fins pratiques, le plaisir qu'ils y prenaient n'existait pas, puisque personne ne le voyait ni ne pouvait le prouver. En quoi ils se trompaient : car Dieu le voyait mieux qu'ils ne le voyaient eux-mêmes. Et Dieu n'a besoin d'aucune preuve pour conclure et pour agir.

7. Sens de l'Enterrement (3) - ou les coups bas

Je crois entrevoir une deuxième raison pourquoi Dieu a voulu que ma propre infidélité durant vingt ans soit suivie d'une infidélité (d'une durée égale - le hasard fait bien les choses...) du plus doué et du plus prestigieux parmi mes élèves. Comme je l'explique longuement dans Récoltes et Semailles, la relation entre cet élève et moi fut une sorte d'affinité élective, une ``lune de miel'' intellectuelle, une communion parfaite dans une passion de connaissance commune, d'une qualité si rare et (surtout) si haute, que l'idée qu'elle pourrait se dégrader jamais m'aurait en ces temps-là paru insensée et quasiment sacrilège ! Or, avec le recul et par un examen patient et obstiné, j'ai pu me rendre compte que dès ces temps idylliques d'une amitié qui semblait sans tâche ni fêlure, le ver déjà était dans le fruit, tant en moi, l'aîné, qu'en lui. En lui comme en moi, ce ``ver'' provenait des pulsions de l'égo qui, à l'abri du regard, sont à l'affut sans cesse pour faire feu de tout bois. Tant que les regards de l'un et de l'autre restaient absorbés par la fascination des ``choses de l'esprit'', ces pulsions cachées pouvaient en prendre à leur aise, pour se déclencher en actes à la première bonne occasion venue.

Mais, dans une relation à autrui, faire attention aux pulsions secrètes de la psyché que celle-ci met en jeu, tant en soi-même (en tout premier lieu) qu'en l'autre, est une \textit{vigilance} de l'âme, d'une essence plus subtile que les dispositions suscitées par la seule pulsion de connaissance intellectuelle (ou artistique, ou amóureuse, le cas échéant...).

Ainsi à mes yeux la mésaventure de mon élève et ami avec moi, et la mienne avec lui, illustre dans un cas particulièrement extrême, dans un cas qui peut passer pour ``idéal'', une vérité essentielle et universellement négligée, \textit{concernant toutes les relations humaines : c'est\unsure{,} qu'en l'absence d'une telle vigilance, d'essence spirituelle, qui naît d'une maturité et qui la promeut à son tour, les relations en apparence même les plus sûres ou les plus anodines, qu'elles soient basées sur des liens intellectuels} (comme ce fut le cas avec mon jeune ami), \textit{ou affectif, ou amoureux, ou de toute autre nature quels qu'ils soient, sont une source infinie de conflits sans issue.} Il en est ainsi, en tous cas, dès qu'une relation est assez forte pour (hem, hem...) pouvoir alimenter une envie de faire des vacheries ou de faire servir l'autre à ses fins, et pendant un temps suffisant pour que (hem, hem...) de telles envies aient le loisir au moins de se concrétiser en actes juteux et substantiels !

Ainsi cette ``mésaventure'' rocambolesque, débarassée de son grinçant par un ``happy end'' édifiant que Dieu tient en réserve pour très bientôt, m'apparaît comme une semonce magistrale de Dieu pour inculquer aux hommes cette ``vérité essentielle'' que je viens de dire et par là-même le sens très concret de la recommandation socratique (souvent citée, et très rarement comprise) du ``\textit{connais-toi toi-même}''.

% archive sheet number: 328
